
\section{Experiência Profissional}

\subsection*{\textbf{NextCompany} \hfill Brasil}
\textit{Desenvolvedor de Software Júnior \hfill Jan 2025 -- Atual}
\begin{itemize}
	\item Atuei no desenvolvimento e evolução do CRM corporativo (wCRM), implantado em nuvem com Docker, participando da transição para o wCompany 2.0, projetado com arquitetura de microsserviços, onde cada rotina do sistema foi implementada como um serviço independente, garantindo escalabilidade, isolamento de responsabilidades e facilidade de manutenção.
	\item Desenvolvi e mantive microsserviços backend utilizando NestJS e Java (Spring Boot), estruturados com princípios SOLID e Clean Architecture, integrados via APIs REST e comunicação assíncrona com Kafka, além do uso de Redis para cache, otimização de desempenho e redução de carga em operações frequentes.
	\item Colaborei em equipes ágeis (Scrum e Kanban) no desenvolvimento de frontend com TypeScript, React e Angular, além da modelagem de dados e criação de consultas SQL complexas em PostgreSQL com Prisma ORM, aplicadas a cadastros, relatórios analíticos, financeiros e gerenciais, incluindo integração com MongoDB para armazenamento de dados fiscais.
\end{itemize}


\subsection*{\textbf{Comunicare Solutions} \hfill Brasil}
\textit{Desenvolvedor Mobile e Backend \hfill Jan 2021 -- Fev 2022}
\begin{itemize}
	\item Desenvolvi aplicação mobile multiplataforma em Flutter voltada ao apoio terapêutico de pessoas com apraxia da fala, com foco em acessibilidade, usabilidade e validação funcional junto a profissionais da área da saúde.
	\item Implementei backend em Python para integração de dados, processamento de informações e comunicação entre sistemas, garantindo confiabilidade e organização do fluxo de dados.
	\item Integrei modelos de reconhecimento de voz utilizando Kaldi em ambiente Linux (Ubuntu), realizando ajustes e execuções shell-based, além de apoiar a análise vocal por meio da geração de espectrogramas em colaboração com fonoaudiólogos.
\end{itemize}
